%preamble document set-up

%-----------------------------------------------------

%Package Includes

%\usepackage{lipsum}
%Generates blind text
\usepackage{graphicx}
%Inclusion of graphics, including resizing etc.
\usepackage[top=2.5cm,bottom=2cm,right=2cm]{geometry}
%For page layout, etc.
\usepackage{amssymb}
%Symbols
\usepackage{setspace} 
%Set linespacing
\usepackage[hidelinks]{hyperref} 
%Clickable links, i.e. \ref
\usepackage[textsize=small]{todonotes} 
%For making to-do notes in the pdf
\usepackage{float} 
%Redefine the float, now we can use 'H' instead of 'h'
\usepackage{array} 
%Extend the array environment
\usepackage{pdfpages} 
%Allows the inclusion of external pdf's 
\usepackage{color, colortbl}
%Allows us to manage the colour of items in the final output 
\usepackage{fancyhdr}
%Allows control of the headers and footers
\usepackage{nameref}
%Refer to sections by their name
\usepackage{chemfig} 
%Draw chemical diagrams
\usepackage{listings} 
%Typeset programming languages
\lstset{language=bash, numbers=left, frame=single, breaklines=true, numberstyle=\footnotesize}
%Set-up date for listings package
\usepackage{framed} 
%Framed or shaded regions that can break across pages
\usepackage[os=win]{menukeys}
%Format menu sequences, paths and keystrokes from lists
\usepackage{draftwatermark} 
%Put a grey textual watermark on document pages
\usepackage{tikz} 
%Used for drawing within the package
\usetikzlibrary{positioning,decorations.markings,patterns,shapes}
%Librarys for use with tikz
\usepackage{multirow} 
%Create tabular cells spanning multiple rows
\usepackage[per-mode=symbol]{siunitx} 
%For putting certain scientific notation and units i.e. x10^.
\usepackage{titlesec} 
%Custom titles
\usepackage[normalem]{ulem}
%Strikethrough text with \sout{} (Not in Maths mode)
\usepackage{cancel}
%Cancel text within maths mode by using \cancel{}
\usepackage{lmodern}
%Use additional symbols such as Numero
\usepackage{textcomp}
%Acronyms package for the control of acronyms
\usepackage[toc,nonumberlist,acronym]{glossaries}

\makeglossaries

\newtheorem{derivation}{Derivation}

%-----------------------------------------------------
%Paragraph and Chapter Styling,

%These are the two commands that allow subsubsections to have numbering and show in the T.O.C.
%The latex numbering system is as follows:

% -1 Part
% 0 Chapter
% 1 Section
% 2 Subsection
% 3 Subsubsection
% 4 Paragraph
% 5 Subparagraph

\setcounter{tocdepth}{2}
\setcounter{secnumdepth}{2}

\titleformat*{\section}{\LARGE\bfseries}
\titleformat*{\subsection}{\Large\bfseries}
\titleformat*{\subsubsection}{\large\bfseries}

%-----------------------------------------------------

%Page and Header Styling
\setlength{\headheight}{14pt}
\pagestyle{fancy}
\renewcommand{\headrulewidth}{0.25pt}
\fancyhead[L]{\textbf{Building Services}}
\fancyhead[R]{\textbf{Contents}}
\setlength\voffset{0in}

%-----------------------------------------------------

%Watermark
\SetWatermarkText{DRAFT}
\SetWatermarkScale{1}
\SetWatermarkLightness{0.9}
%\SetWatermarkColor{red}

%-----------------------------------------------------

%Document meta-data

\hypersetup{
    pdftitle={Calculation Methods for Building Services System Design},
    pdfauthor={Adam Rees},
    pdfsubject={Building Services},
    pdfcreator={Adam Rees},
	pdfstartview={XYZ null null 1.00},
}
    
%-----------------------------------------------------

%Custom Commands

\definecolor{Grey}{gray}{0.95}
\lstset{language=bash, numbers=left, frame=single, breaklines=true, numberstyle=\footnotesize}

\DeclareSIUnit\lsp{l/s/p}

%-----------------------------------------------------

%Include Glossary

%\newacronym{LABEL}{ABBREVIATION}{FULL}

\newacronym{ACH}{ACH}{Air Changes per Hour}
\newacronym{CWS}{CWS}{Cold Water Service}
\newacronym{HWS}{HWS}{Hot Water Service}
\newacronym{VOC}{VOC}{Volatile Organic Compounds}
\newacronym{AHU}{AHU}{Air Handling Unit}
\newacronym{HTM}{HTM}{Health Technical Memorandum}
\newacronym{HBN}{HBN}{Health Building Note}
\newacronym{BB101}{BB101}{Building Bulletin 101 - Guidelines on ventilation thermal comfort and indoor air quality in schools}
\newacronym{CIBSE}{CIBSE}{Chartered Institution of Building Services Engineers}
\newacronym{ASHRAE}{ASHRAE}{American Society of Heating, Refrigerating and Air-Conditioning Engineers}
\newacronym{BSI}{BSI}{British Standards Institution}
\newacronym{BS}{BS}{British Standard}
\newacronym{EN}{EN}{European Norm}
\newacronym{ISO}{ISO}{International Organization for Standardization}